\section{Conclusion}
\label{sec:conclusion}

This paper presented a dual-stream Kalman transformer architecture for wearable fall detection that achieves state-of-the-art performance on the SmartFallMM dataset. Our comprehensive experimental evaluation, conducted using rigorous 21-fold Leave-One-Subject-Out Cross-Validation (LOSO-CV), demonstrates several key contributions:

\textbf{1. Kalman fusion enables effective gyroscope utilization.} While prior work found that adding raw gyroscope data to accelerometer-only models degraded performance, we demonstrate that Kalman-fused orientation features (roll, pitch, yaw) improve transformer performance by +1.40\% F1 (89.49\% $\rightarrow$ 90.89\%). The Kalman filter transforms noisy angular velocity into semantically meaningful body pose information that the transformer's attention mechanism can effectively leverage.

\textbf{2. Dual-stream architecture with channel-aware normalization.} Processing accelerometer and orientation streams separately before fusion, combined with normalizing only accelerometer channels (keeping orientation in raw radians), achieves our best result of 91.10\% F1 score. This represents a +6.36\% improvement over the accelerometer-only baseline (84.49\%).

\textbf{3. Transformer outperforms CNN.} Under identical training conditions, the dual-stream Kalman transformer (91.10\% F1) outperforms the CNN baseline (88.86\% F1) by +2.23\%. The Transformer wins on 14/21 subjects (66.7\%), with particular advantage on challenging subjects exhibiting atypical fall patterns.

\textbf{4. Architecture-dependent Kalman effects.} While transformers benefit from Kalman fusion, CNNs perform better with raw input. This suggests that CNNs rely on high-frequency local patterns (impact spikes) that Kalman filtering smooths away, whereas transformers leverage the global semantic information in orientation trajectories.

\textbf{5. Failure mode characterization.} Subjects with F1 $<$ 80\% exhibit 20\% lower impact forces during falls, making their patterns harder to distinguish from ADLs. This represents a fundamental challenge that future work should address through data augmentation, subject-specific calibration, or multi-modal fusion.

The 91.10\% F1 score achieved by our dual-stream Kalman transformer represents the highest reported performance on the SmartFallMM dataset using wearable IMU data, demonstrating the effectiveness of our approach for real-world fall detection applications on resource-constrained wearable devices.
